% ============================================================
% Attacker Cognitive Loop: An Interpretable Rule-Based Engine
% for Estimating, Predicting, and Engaging Adaptive Attackers
% arXiv submission (cs.CR), draft v0.2
%
% Compile: pdflatex (2 passes). Only standard packages are used
% so the source builds on arXiv.
% ============================================================
\documentclass[11pt]{article}

\usepackage[margin=1in]{geometry}
\usepackage{amsmath,amssymb}
\usepackage{booktabs}
\usepackage{adjustbox}
\usepackage{url}
\usepackage[hidelinks]{hyperref}
\usepackage{xcolor}
\usepackage{verbatim}
\setlength{\emergencystretch}{2em}
\raggedbottom

\newcommand{\act}[1]{\texttt{#1}}
\newcommand{\intent}[1]{\textsc{#1}}

\title{\bfseries Attacker Cognitive Loop: An Interpretable Rule-Based Engine for Estimating, Predicting, and Engaging Adaptive Attackers}

\author{Jiahao Lai \\
  \small ASSCOR Project \\
  \small \texttt{chins-xing@proton.me}}

\date{\today}

\begin{document}

\maketitle

\begin{abstract}
Attackers and defenders operate under fundamentally unequal information endowments: attackers choose when, where, and how to probe; defenders react to signals that arrive too late and too sparse. This paper presents the Attacker Cognitive Loop (ACL), a system designed to actively reverse this information asymmetry --- not by blocking attacker actions, but by guiding them through lightweight decoys that reveal intent earlier than the attacker intends. ACL closes the loop between observation and intervention, turning the defender's information disadvantage into a controllable orchestration problem. The engine is implemented as an interpretable, rule-based extension of the open-source ASSCOR platform, with every component mapped to a public repository path and branch. Its heuristic weights are explicitly not optimal; they are provided solely to demonstrate closed-loop mechanics and to serve as a baseline for community-driven replacement. We formulate the research question --- \textit{how AI, shared tooling, and individual experience jointly shape the attacker's action distribution, and whether defenders can actively reshape it through observation, deception, and guidance} --- and state ten open research questions (Q1--Q10) with evidence ratings. We provide the complete formal specification of the scoring, normalization, sharpness, and utility functions, a reproducible four-round closed-loop trace with exact numerical values, a sensitivity analysis of the hand-tuned weights, and simulated per-round latency benchmarks.
\end{abstract}

\section{Introduction}
\label{sec:intro}

Classical defense embeds blocking controls: access control and isolation make
undesirable states unreachable. Surveys of moving-target defense (MTD) have
documented two inherent weaknesses of this paradigm. First, a
\textit{cognitive limitation}: defenders cannot enumerate all latent attacker
paths, and attackers enjoy an information advantage during prolonged
reconnaissance. Second, a \textit{temporal asymmetry}: attackers can plan
over long horizons while defenders must respond in real time, and patch lag
leaves a window of opportunity.

At the heart of the two weaknesses above lies a single structural fact: attackers hold an inherent information advantage. They observe the defender's environment at their chosen pace, select which paths to probe, and can wait for the defender's response patterns to emerge. Defenders, by contrast, see only the artifacts of actions already taken --- alerts, logs, and blocked connections. This is not a technical failure of firewalls or IPS; it is a cognitive asymmetry built into the defender's observational posture.

Guided active defense, as we define it, is an attempt to reverse this asymmetry. Instead of waiting for the attacker to reveal intent through damage, the defender proactively deploys ephemeral decoys that appear to reduce the attacker's uncertainty --- inducing the attacker to verify hypotheses early, thereby exposing TTPs and intent before the attack reaches its objective. ACL is designed as an orchestration engine for exploring this reversal: it estimates what the attacker believes, predicts what the attacker will do next, and selects decoys that maximize information gain while minimizing the defender's exposure. The current implementation provides a fully auditable rule-based baseline against which more principled models can be compared.

Guided active defense does not aim to block everything; it aims to create an
\textit{information asymmetry} in the defender's favor. By deploying
lightweight decoy targets and fabricated intelligence, the defender induces
the attacker to abandon slow, patient operation in favor of rapid
verification --- exposing tactics, techniques, and procedures (TTPs) and
intent earlier than they would otherwise be revealed. Intelligence collected
from each interaction feeds the next, more accurate round of deception,
forming a positive feedback loop: \textit{guide $\rightarrow$ collect
$\rightarrow$ analyze $\rightarrow$ decide $\rightarrow$ guide again}.

In the MITRE Engage taxonomy, engagement is used here as a guidance tool,
not a blocking tool; host isolation is demoted to a fallback when guidance
fails or confirms high risk.

This work makes the following contributions:
\begin{enumerate}
\item We formulate the attacker-cognitive-loop research question that
      connects attacker modeling with adaptive deception.
\item We present an interpretable, rule-based engine (ACL) that realizes the
      loop: state estimation, action-distribution prediction, utility-driven
      engagement selection, and sharpness-driven strategy switching, with
      the complete formal specification of every score, normalization,
      sharpness, and utility term (\S\ref{sec:prediction-formal},
      \S\ref{sec:sharpness-formal}, \S\ref{sec:utility-formal}).
\item We provide a reproducible four-round closed-loop trace with exact
      numerical values computed from the reference implementation's weights
      (\S\ref{sec:trace}), so the engine's behavior can be independently
      verified and replayed.
\item We anchor the work in a real, public codebase: the engine is an
      implemented extension of the open-source ASSCOR assessment platform,
      with every component mapped to a repository path and branch
      (\S\ref{sec:repo}, \S\ref{tab:repo-map}), so the paper is directly
      auditable rather than a paper-only design.
\item We state ten open research questions (Q1--Q10) about attacker
      behavior in the AI era, each with an evidence rating, and a
      validation plan.
\item We are transparent about what is implemented (the engine) versus what
      remains to be answered (the questions), and we analyze the
      sensitivity of the hand-tuned weights (\S\ref{sec:sensitivity}). No
      assumption is presented as a conclusion.
\end{enumerate}

\section{Related Work}
\label{sec:related}

\textbf{Decision-theoretic planning.} Shinde and Doshi's
I-POMDP\textsuperscript{2} \cite{shinde-doshi} extends partially observable
Markov decision processes to model the defender's beliefs about the
attacker's beliefs, capabilities, and preferences, with finite nested
recursive reasoning. Their experiments observe level-3 defenders who first
mislead the attacker into believing the defense is passive, then deploy
decoys --- a decision-theoretic justification for guidance over blocking. A
full I-POMDP\textsuperscript{2} solver is computationally prohibitive; our
engine is a transparent rule-based approximation that can be replaced by
richer models without changing the loop architecture.

\textbf{Game-theoretic deception.} Zhu's tutorial \cite{zhu} and Anwar \&
Kamhoua's attack-graph games \cite{anwar-kamhoua} formalize signaling games,
dynamic games, mechanism design, and hypergames for deception. General-sum
stochastic games make pure-strategy Nash equilibrium computation
PSPACE-hard, and equilibrium concepts assume symmetric optimality that
poorly models the attacker's \textit{bounded rationality}. We therefore
prefer the decision-theoretic view.

\textbf{Attacker modeling and attribution.} The Diamond Model of intrusion
analysis \cite{diamond} structures every intrusion event as four vertices
--- \textit{adversary}, \textit{capability}, \textit{infrastructure}, and
\textit{victim} --- connected by six phases of an activity thread, with
meta-features (timestamp, kill-chain phase, result, direction, method,
resources) describing each event. Its central axiom, that each intrusion is
an event in an activity thread whose evolution is governed by the adversary,
is the analytical backbone we reuse: our six-layer graph
(\S\ref{sec:graph}) is the Diamond Model's relational projection onto a
temporal substrate --- the Adversary layer holds the adversary vertex, the
Capability layer holds capabilities and TTPs, the Network/Dependency layers
hold infrastructure, and the Identity/Evidence layers hold the victim and
the meta-features. Unlike the Diamond Model, which is a retrospective
analytic framework, ACL consumes the same relational structure
\textit{forward} to estimate the attacker's state and predict its next
action. We also note the fundamental attribution problem raised by the
shared-capability questions (Q5/Q6, \S\ref{sec:hypotheses}): observed
similarity between intrusions does not imply the same adversary, since
adversaries share tooling, code, and (increasingly) AI-generated strategy.

\textbf{Defensive deception and honeypot limitations.} Zhang and Thing
\cite{zhang-thing} survey three decades of honeypots, honeytokens, and MTD
across a four-layer stack (network/system/software/data) mapped to the kill
chain, and document the two structural weaknesses of honeypot-only defenses
that motivate our guided-engagement design: (i) \textit{static deployment}:
if honeypots and honeytokens are left with static configuration, an
adversary has enough time to infer their existence, map them, and evade
them; and (ii) \textit{pivot risk}: high-interaction honeypots, which offer
a real operating-system environment, can themselves be compromised and used
as a pivot point to attack other systems. Two further weaknesses are
well-known in practice: (iii) \textit{low-interaction honeypots have low
fidelity}, so sophisticated attackers can fingerprint and bypass them; and
(iv) \textit{passivity}: traditional honeypots wait for the attacker and
yield little intelligence about intent until the attacker has already
committed. Our loop addresses (i) by making engagement a \textit{function of
the predicted action distribution} that changes each round, (ii) by the
\textit{sufficiency principle} (\S\ref{sec:decoy-mapping}) --- decoys expose
only a plausible break-in appearance with no real system to compromise ---
(iii) by keeping decoys lightweight and cheap enough to deploy at
single-host scale, and (iv) by using decoys as \textit{active sensors} whose
interaction chains feed the state engine.

\textbf{MTD design principles.} MTD surveys \cite{mtd-survey} contribute
five design principles (coverage, unpredictability, timeliness,
superstability, functional equivalence) that constrain any future MTD
component. We adopt honeypot/honeytoken as the primary guidance mechanism
and treat MTD as an optional deterministic-weakening component.

\textbf{Behavior standardization.} MITRE ATT\&CK \cite{attck} provides the
semantic layer for attack behavior. Our prediction engine maps action
distributions onto TTP distributions rather than replacing ATT\&CK, and our
state engine infers intent from observed TTPs through an explicit
technique$\to$intent table. Table~\ref{tab:ttp-intent} lists the mapping
used by the reference implementation (the full table is
\S\ref{sec:state-engine}); the same table drives the reverse projection
$P(\mathit{TTP}_i) = \sum_j P(\mathit{TTP}_i \mid \mathit{Action}_j)\,
P(\mathit{Action}_j)$ (\S\ref{sec:prediction-formal}). This makes the
ATT\&CK$\to$state-machine bridge explicit, auditable, and replaceable:
unknown techniques remain \act{unknown} rather than being guessed.

\begin{table}[ht]
\centering
\begin{adjustbox}{max width=\linewidth}\begin{tabular}{@{}lll@{}}
\toprule
\textbf{ATT\&CK technique} & \textbf{Inferred intent} & \textbf{Example} \\
\midrule
T1595, T1592, T1590, T1046 & recon & Active Scanning (T1595) \\
T1110, T1003, T1555, T1078 & credential & Brute Force (T1110) \\
T1021, T1570, T1550, T1028 & lateral & Remote Services (T1021) \\
T1567, T1048, T1030, T1537 & data\_theft & Exfiltration Over Web Service (T1567) \\
T1190, T1193, T1566 & web\_attack & Exploit Public-Facing Application (T1190) \\
\bottomrule
\end{tabular}\end{adjustbox}
\caption{TTP $\rightarrow$ intent mapping used by the state engine
(\S\ref{sec:state-engine}). Techniques outside this table remain
\act{unknown}; the mapping is open data, not engine logic.}
\label{tab:ttp-intent}
\end{table}

\textbf{Risk propagation.} SRD (Systemic Risk Dynamics) models risk
diffusion over a network graph \cite{asscor-wp}; we reuse its substrate for
the topology layer and as the dynamic state-evolution layer of the
cognitive loop.

Our work is built on ASSCOR (ASSess + CORE), an open-source security acceptability assessment platform. ASSCOR provides the three substrates that ACL requires: (i) a six-layer temporal graph (Network, Dependency, Identity, Attacker, Capability, Evidence) that carries observations and alerts as temporal events; (ii) SSAM 2.0, a multidimensional computable assessment model that produces per-host acceptability scores $\sigma \in [0,100]$; and (iii) Prism/SRD, a three-layer risk propagation engine that models how compromise diffuses through the network. Without these substrates, ACL's cognitive loop would have no data to consume --- the state engine needs graph topology to localize observations, the predictor needs SSAM scores as vulnerability signals, and the engagement planner needs Prism risk edges to prioritize containment targets.

ACL is thus not a replacement for these engines; it is an orchestration layer that closes the loop between assessment (passive) and intervention (active). ASSCOR's main branch provides the stable assessment baseline; ACL lives in the ASSCOR-Research-Core branch, an experimental sandbox explicitly designated for active-defense prototyping. This split is not merely a code-management convenience --- it reflects the different development rhythms of assessment (stable, auditable) and active defense (iterative, replaceable), and ensures that experimental prediction logic never compromises the production-grade assessment pipeline.

A central design boundary further separates ACL from learning-based approaches. Neural-network-based predictors have demonstrated impressive performance in low-stakes security tasks such as log correlation, anomaly scoring, and alert triage. However, ACL operates in a regime where individual recommendations carry operational risk: a decoy deployed on the wrong target, or at the wrong time, can alert the attacker, waste defensive resources, or be pivoted against. In this regime, interpretability under failure is a safety requirement, not a desirable extra. When a defensive intervention fails, the operator must be able to trace the decision to specific, auditable factors --- not to aggregate model weights. ACL therefore defaults to a fully decomposable rule-based scorer (Eq.~\ref{eq:score}). The architecture remains open to neural replacements at Seam 2 (\S\ref{sec:replacement}), with the caveat that any proposed replacement must preserve a mechanism for post-hoc explanation (e.g., through surrogate models or attention-based attribution) so that the safety requirement is not violated.

\section{Problem Formulation}
\label{sec:problem}

\subsection{Attacker State}
\label{sec:state}

We model the attacker's cognitive state as:
\[
\begin{aligned}
\mathit{AttackerState} = (&\mathit{Capability},\mathit{Experience},\mathit{Intent},
\mathit{TTP\_Repertoire},\mathit{SharedCapability},\\
&\mathit{IndividualCapability},\mathit{AI\_Dependence},\mathit{TargetKnowledge},
\mathit{BeliefState},\mathit{Objective})
\end{aligned}
\]
where $\mathit{Intent} \in \{\intent{recon}, \intent{credential},
\intent{lateral}, \intent{data\_theft}, \intent{web\_attack}\}$,
$\mathit{AI\_Dependence} \in [0,1]$ measures reliance on AI for
decision/execution, and $\mathit{BeliefState}$ captures the attacker's
subjective perception of the target environment.

\subsection{Prediction Target}
\label{sec:prediction-target}

Rather than a deterministic next-TTP, we predict:
\[
P(A_{t+1} \mid \mathit{AttackerState}_t, \mathit{TargetState}_t,
\mathit{Observation}_t)
\]
and additionally output \textit{most likely}, \textit{most dangerous}, and
\textit{most observable} actions (and a ranked list), feeding both risk
prioritization and engagement selection.

\subsection{The Cognitive Loop}
\label{sec:loop}

The closed loop is:
\[
\begin{aligned}
\mathit{Observation} &\to \mathit{State\ Estimation}
\to \mathit{Behavior\ Prediction}
\to \mathit{Guidance/Deception} \\
&\to \mathit{Attacker\ Response}
\to \mathit{New\ Evidence}
\to \mathit{State\ Update}
\to \mathit{Predict\ Again}
\end{aligned}
\]

\section{System Design: The Attacker Cognitive Loop}
\label{sec:design}

ACL is implemented as a modular, pure-function library (no external
dependencies per component), designed so each component can be independently
tested and replaced. It follows a microkernel architecture: the core kernel
remains untouched; these components are optional, build-tag modules.

\subsection{Multi-Layer Temporal Graph}
\label{sec:graph}

A single \act{Host--connects--Host} graph cannot express attackers,
identities, tools, code, TTPs, and evidence. We maintain six semantic layers
--- Network, Dependency, Identity, Attacker, Capability, Evidence --- over a
shared temporal graph: nodes and edges carry
$\mathit{FirstSeen}/\mathit{LastSeen}$ timestamps, edges carry lifecycle
status (up/down) and attributes (latency, bandwidth, ACL), and nodes carry an
event stream (observations, alerts, decoy triggers) with confidence.
Reachability queries (shortest path, bounded-hop traversal) support
multi-hop analysis. The layer/type/edge-type vocabularies are open strings:
predefined constants cover the common semantics, and any model can introduce
new values without engine changes.

\subsection{Attacker State Engine}
\label{sec:state-engine}

The state engine implements
$\mathit{OldState} + \mathit{Evidence} \to \mathit{NewState}$ with explicit,
interpretable rules. Formally, let the state at round $t$ be
$s_t = (K_t, E_t, I_t, \mathcal{R}_t, \mathcal{C}_t, B_t)$ where $K_t$ is
target knowledge, $E_t$ experience, $I_t$ intent, $\mathcal{R}_t$ the TTP
repertoire, $\mathcal{C}_t$ the capability set, and $B_t$ the belief map.
An evidence item $e = (\mathit{ttp}, \mathit{intent}, \mathit{target},
\mathit{outcome}, c)$ with confidence $c \in [0,1]$ updates the state by:
\begin{align}
K_{t+1} &= \min(1,\; K_t + 0.1\, c \cdot \mathbf{1}[\mathit{target}
\neq \emptyset]) \label{eq:know} \\
E_{t+1} &= \min(1,\; \max(0,\; E_t + 0.05\,\mathbf{1}[\mathit{outcome}
= \text{success}] - 0.03\,\mathbf{1}[\mathit{outcome} = \text{failure}]))
\label{eq:exp} \\
B_{t+1}[\mathit{target}] &= \min(1,\; B_t[\mathit{target}]
+ 0.15\, c) \label{eq:belief} \\
\mathcal{R}_{t+1} &= \mathcal{R}_t \cup \{\mathit{ttp}\},
\quad \mathcal{C}_{t+1} = \mathcal{C}_t \cup \{\mathit{cap}(\mathit{ttp})\}
\label{eq:repertoire} \\
I_{t+1} &= \arg\max_{e'} c_{e'} \quad \text{over evidence with known intent,} \label{eq:intent} \\
&\qquad \text{where explicit $\mathit{intent}$ overrides $\mathit{IntentFromTTP}(\mathit{ttp})$} \nonumber
\end{align}
where $\mathit{cap}(\cdot)$ maps a technique to its capability requirement.
Updates are pure: the input state is never mutated. Intent inference is
deliberately conservative --- unknown TTPs remain \act{unknown} rather than
being guessed (see Algorithm~1 in \S\ref{sec:impl-update}).

\subsection{Prediction Engine}
\label{sec:predictor}

\subsubsection{Raw Scoring}
\label{sec:prediction-formal}

The prediction engine converts state plus target state into a normalized
action distribution over six actions $a \in \mathcal{A} = \{\act{recon},
\act{credential}, \act{lateral}, \act{data\_theft}, \act{web\_attack},
\act{maintain}\}$. Let the state at round $t$ be $s_t$ with intent $I_t$,
experience $E_t$, target knowledge $K_t$, and AI dependence $A_t$; and let
the target have SSAM score $\sigma \in [0,100]$ and exposure
$\rho \in [0,1]$. The raw score of action $a$ is the sum of six
interpretable terms:
\begin{equation}
\label{eq:score}
\begin{aligned}
\mathrm{score}_t(a) ={}& b_a + w_{\mathrm{int}}\mathbf{1}[I(a)=I_t]
+ E_t\mathbf{1}[a\in\mathcal{A}_{\mathrm{exec}}]
-0.5E_t\mathbf{1}[a=\act{recon}]\\
&+ K_t\mathbf{1}[a\in\mathcal{A}_{\mathrm{kn}}]
+ A_t\mathbf{1}[a\in\mathcal{A}_{\mathrm{ai}}]
-0.5A_t\mathbf{1}[a=\act{recon}]\\
&+ w_{\mathrm{vuln}}v_t\mathbf{1}[a\in\mathcal{A}_{\mathrm{vuln}}].
\end{aligned}
\end{equation}
with
\begin{equation}
v_t \;=\; \frac{100-\sigma}{100} + \rho
\label{eq:vuln}
\end{equation}
and the action subsets are
\[
\begin{aligned}
\mathcal{A}_{\mathrm{exec}} &= \{\act{lateral}, \act{data\_theft}\}, &
\mathcal{A}_{\mathrm{kn}} &= \{\act{credential}, \act{lateral}, \act{data\_theft}\},\\
\mathcal{A}_{\mathrm{ai}} &= \{\act{credential}, \act{web\_attack}\}, &
\mathcal{A}_{\mathrm{vuln}} &= \{\act{credential}, \act{lateral}, \act{data\_theft}, \act{web\_attack}\}.
\end{aligned}
\]
The baseline priors and weights used by the reference implementation are
listed in Table~\ref{tab:params}. All scores are clamped to $[0,\infty)$
before normalization.

\begin{table}[ht]
\centering
\begin{adjustbox}{max width=\linewidth}\begin{tabular}{@{}ll@{}}
\toprule
\textbf{Parameter} & \textbf{Default value} \\
\midrule
$b_a$ (baseline prior) & recon 1.0, credential 1.2, lateral 1.1, \\
 & data\_theft 0.9, web\_attack 1.0, maintain 0.5 \\
$w_{\mathrm{int}}$ (intent continuity) & 3.0 \\
$w_{\mathrm{vuln}}$ (vulnerability weight) & 0.8 \\
$T$ (softmax temperature) & 1.0 \\
danger weight (MostDangerous) & data\_theft 1.0, lateral 0.8, credential 0.6, \\
 & web\_attack 0.5, recon 0.3, maintain 0.2 \\
$\alpha,\beta,\gamma,\delta$ (utility) & 1.0, 0.8, 0.6, 1.2 \\
sharpness threshold $\theta$ & 0.3 \\
\bottomrule
\end{tabular}\end{adjustbox}
\caption{Default weights and priors of the reference implementation
(Table~\ref{eq:score}, \ref{eq:utility}). All values are configuration, not
engine logic.}
\label{tab:params}
\end{table}

\subsubsection{Normalization}
\label{sec:softmax}

Scores are converted to a probability distribution by a temperature-
controlled softmax (numerically stabilized by subtracting the maximum
score):
\begin{equation}
\label{eq:softmax}
P_t(a) \;=\; \frac{\exp\!\big((\mathrm{score}_t(a) -
\max_{a'}\mathrm{score}_t(a'))/T\big)}
{\sum_{a'' \in \mathcal{A}} \exp\!\big((\mathrm{score}_t(a'') -
\max_{a'''}\mathrm{score}_t(a'''))/T\big)}
\end{equation}
which satisfies $P_t(a) \geq 0$ and $\sum_{a \in \mathcal{A}} P_t(a) = 1$.
The distribution is then projected onto the ATT\&CK TTP layer:
\begin{equation}
\label{eq:ttp-proj}
P(\mathit{TTP}_i) \;=\; \sum_{a \in \mathcal{A}}
P(\mathit{TTP}_i \mid a)\; P_t(a),
\qquad
P(\mathit{TTP}_i \mid a) = 1/|\mathcal{T}_a| \;\text{for } \mathit{TTP}_i
\in \mathcal{T}_a,
\end{equation}
where $\mathcal{T}_a$ is the technique set associated with action $a$
(Table~\ref{tab:ttp-intent} and its reverse). This keeps ATT\&CK as the
behavioral semantic layer rather than replacing it.

\subsubsection{Multi-Outputs}

Beyond the full distribution, the engine outputs:
\begin{align}
a^{\mathrm{ML}}_t &= \arg\max_a P_t(a) \label{eq:most-likely} \\
a^{\mathrm{MD}}_t &= \arg\max_a P_t(a)\, d_a \quad\text{(danger-weighted)}
\label{eq:most-dangerous} \\
a^{\mathrm{MO}}_t &= \arg\max_{a \in \{\act{recon},\act{credential},
\act{lateral}\}} P_t(a) \quad\text{(most observable)}
\label{eq:most-observable}
\end{align}
with danger weights $d_a$ from Table~\ref{tab:params}, feeding risk
prioritization and engagement selection.

\subsection{Engagement Planner}
\label{sec:engagement}

\subsubsection{Utility Function}
\label{sec:utility-formal}

The engagement planner selects deceptive interventions by an explicit
utility function. For an intervention $E$ targeting action $a_E$ with
deployment cost $\mathit{cost}_E$ and exposure $\mathit{exp}_E$, the utility
is:
\begin{equation}
\label{eq:utility}
\mathrm{Utility}(E) \;=\;
\alpha\,\underbrace{P_t(a_E)\; \kappa(\sigma,\rho)}_{\mathrm{IG}(E)}
\;+\; \beta\,\underbrace{\mathit{DP}_E}_{\text{detection probability}}
\;+\; \gamma\,\underbrace{\mathit{AV}_E}_{\text{attribution value}}
\;-\; \delta\,\underbrace{\mathit{cost}_E \cdot \mathit{exp}_E}_{\mathrm{Risk}(E)}
\end{equation}
with the target-coverage factor
\begin{equation}
\label{eq:coverage}
\kappa(\sigma,\rho) \;=\; 0.5 + 0.5 \cdot \frac{(100-\sigma)/100 + \rho}{2}.
\end{equation}
Information gain $\mathrm{IG}(E)$ is the predicted probability of the action
the decoy targets, scaled by coverage; $\mathit{DP}_E$ and $\mathit{AV}_E$
are intrinsic decoy capabilities (Table~\ref{tab:decoy}); Risk is the
product of cost and exposure. Only interventions with positive utility are
candidates; they are returned in descending utility order.

\begin{table}[ht]
\centering
\begin{adjustbox}{max width=\linewidth}\begin{tabular}{@{}llccc@{}}
\toprule
\textbf{Decoy} & \textbf{Target action} & \textbf{DP} & \textbf{AV} & \textbf{Cost/Exposure} \\
\midrule
fake SSH & lateral & 0.90 & 0.60 & 0.30 / 0.20 \\
fake credential & credential & 0.80 & 0.90 & 0.20 / 0.15 \\
fake document & data\_theft & 0.75 & 0.85 & 0.20 / 0.15 \\
fake web & web\_attack & 0.70 & 0.50 & 0.40 / 0.30 \\
scan port & recon & 0.85 & 0.30 & 0.10 / 0.10 \\
\bottomrule
\end{tabular}\end{adjustbox}
\caption{Default decoy catalog: detection probability (DP), attribution
value (AV), and cost/exposure per decoy type.}
\label{tab:decoy}
\end{table}

\subsubsection{Decoy as Sensor}
\label{sec:decoy-mapping}

Following the \textit{sufficiency principle}, decoys need only create a
plausible break-in appearance, not high-fidelity honeypot credentials.
Decoys are treated as sensors: an interaction chain (connect $\to$
credential attempt $\to$ command $\to$ system discovery $\to$ follow-up) is
recorded as a $\mathit{DeceptionRecord}$ and converted into an
$\mathit{Evidence}$ item (with the TTP/intent inferred from the interaction
chain, conservatively: only when a TTP is identified) that feeds the state
engine, closing the loop (\S\ref{sec:loop}).

\subsubsection{Decoy Anti-Pivot and Self-Destruct}
\label{sec:anti-pivot}

Because the engagement planner operates as a sensor, its decoys must not
become a second attack surface: an adversary who compromises a decoy could
otherwise pivot from it onto production assets. All deployed decoys are
therefore \textbf{ephemeral} ($\mathrm{TTL} \leq 60$\,s). If a decoy
receives any \textit{outbound} connection request, it immediately terminates
the session and resets to a clean state. The decoy environment is
\textbf{air-gapped from production assets}: it holds no real credentials, no
routing tables, and no reachable references to the production network.
Combined with the sufficiency principle (\S\ref{sec:decoy-mapping}), this
guarantees that a decoy can collect at most one bounded interaction chain
before it disappears, and that nothing reachable from a decoy can be used to
reach production. The reference implementation enforces this contract in the
decoy deployer (lifecycle timer + outbound-connection teardown); the
corresponding defensive-fallback strategy (\S\ref{sec:strategy}) further
ensures that containment, not decoy deployment, is the default response to
unknown intent or unknown TTPs.

\subsection{Closed-Loop Controller}
\label{sec:controller}

\subsubsection{Decision Sharpness}
\label{sec:sharpness-formal}

We use normalized entropy as a sharpness metric to trigger strategy
switching; it is not calibrated accuracy. Let
\begin{equation}
\label{eq:sharpness}
S_t \;=\; 1 - \frac{H(P_t)}{\ln |\mathcal{A}|},
\qquad
H(P_t) = -\sum_{a \in \mathcal{A}} P_t(a) \ln P_t(a),
\end{equation}
so sharp distributions (low entropy) have high sharpness and flat ones low.
\textbf{This metric is a heuristic for action-distribution peakedness, not a
probabilistic guarantee of correctness.} A flat distribution --- the honest
expression of an ambiguous attacker --- yields low sharpness, which drives
the controller toward containment or collection rather than aggressive
engagement; the strategy switch (\ref{eq:strategy}) never interprets low
entropy as high certainty.

\subsubsection{Strategy}
\label{sec:strategy}

The controller selects a strategy by first applying a defensive fallback
and then thresholding sharpness:
\begin{equation}
\label{eq:strategy}
\text{strategy}_t \;=\;
\begin{cases}
\act{contain} & I_t = \intent{unknown} \text{ or } \exists\, \text{unknown
TTP in observations} \\
\act{collect} & S_t < \theta \quad\text{(keep only the highest-IG
intervention: intelligence first)} \\
\act{engage} & S_t \geq \theta \quad\text{(keep all positive-utility
interventions: active guidance)}
\end{cases}
\end{equation}
with default $\theta = 0.3$. \textbf{Contain} denotes temporary network
isolation or process suspension of the affected target until new evidence
arrives; no decoys are deployed while containing. It guarantees that the
worst case --- an attacker using a technique absent from the TTP$\to$intent
mapping, or an intent that cannot be established --- leaves the system in a
safe state instead of a passively open one. This realizes the
information-gain closed loop: predict $\to$ select $\to$ observe $\to$
update $\to$ re-predict.

\section{Design Rationale}
\label{sec:rationale}

This section documents the concrete design choices and the trade-offs they
entail, so that the system can be evaluated and modified by others.

\subsection{Rule-Based over Learned or Game-Theoretic Models}
The state engine and predictor use hand-crafted rules and weights. This choice reflects the operational constraints of security decision-making, not a temporary expedient.

First, accountability under failure. In security, the cost of a wrong decision is asymmetric and operationally consequential. A decoy deployed against the wrong attacker intent does not merely reduce accuracy; it may signal the defender's posture to the attacker, or waste the limited window in which the attacker is observable. When such a failure occurs, the operator needs to know why. A neural predictor, by itself, cannot answer that question without post-hoc explanation tools, which introduce their own approximation errors and are not yet standardized for audit trails. ACL's rule-based scorer, by contrast, emits a decomposition of every score term at prediction time (Eq.~\ref{eq:score}), so every decision can be inspected, audited, and, if necessary, corrected. This is a non-negotiable requirement for a system that recommends actions with real operational impact.

Second, no training data. The research questions Q1--Q10 are not yet answered; there is no labeled dataset of attacker behavior from which to learn. Rules are the only honest starting point. But even with abundant data, we would not replace the scoring function with a black-box model without preserving full decomposability --- the safety requirement is structural, not data-dependent.

Third, computational cost. Full I-POMDP\textsuperscript{2} or general-sum stochastic-game solution is computationally prohibitive in real-time constraints. A transparent approximation is acceptable as a first implementation, and the loop architecture is deliberately model-agnostic so that better white-box approximations can replace the current rules without sacrificing auditability.

We emphasize: this is a constitutional principle of the ASSCOR platform, not a preference. ACL is designed to prioritize determinism and traceability over raw predictive accuracy. The three replacement seams (\S\ref{sec:replacement}) are explicitly provided so that the community can substitute more principled models --- but any replacement must preserve explainability, or it violates the safety contract of the system.

\subsection{Intent-Continuity Weight}

The action whose intent equals the current $\mathit{AttackerState.Intent}$
receives a strong additive weight ($w_{\mathrm{int}} = 3.0$, versus baseline
0.5--1.2). Rationale: the current stage is the strongest single predictor of
the next action --- an attacker currently harvesting credentials tends to
continue credential actions. The magnitude was tuned so that, under a clear
intent, the distribution becomes sharp enough for the sharpness-driven
controller to switch to active engagement, while remaining responsive to new
contradictory evidence. The sensitivity of this choice is analyzed in
\S\ref{sec:sensitivity}.

\subsection{Entropy-Based Sharpness}

Decision sharpness is defined as the normalized entropy of the action
distribution (Eq.~\ref{eq:sharpness}): sharp distributions have high
sharpness, flat ones low. It is a \textit{peakedness heuristic}, not
calibrated accuracy --- we state this explicitly in \S\ref{sec:discussion}.
The default threshold ($\theta = 0.3$) separates the collect strategy (flat
distributions: prefer intelligence gathering) from the engage strategy
(sharp distributions: allow active guidance), and the defensive fallback
(\S\ref{sec:strategy}) overrides both whenever intent is unknown or an
unknown TTP is observed. Both the threshold and the temperature are
configuration parameters.

\subsection{Utility Weights}

The engagement utility (Eq.~\ref{eq:utility}) uses default weights
$\alpha=1.0$, $\beta=0.8$, $\gamma=0.6$, $\delta=1.2$. Information gain
dominates because intelligence acquisition is the primary objective of the
loop (\S\ref{sec:loop}); risk carries the largest coefficient because decoy
exposure (detection as fake, or collateral business impact) is the main cost
of deception. These weights are tunable per deployment.

\subsection{Open Semantic Vocabulary}
\label{sec:open-vocab}

$\mathit{Layer}$, $\mathit{NodeType}$, and $\mathit{EdgeType}$ are plain
strings. Predefined constants cover the common semantics (six graph layers,
23 node types, nine edge types), but any model can introduce new values
without changing the engine. This decouples the graph engine from any
specific domain vocabulary --- the same engine serves network topology,
identity, attacker, capability, and evidence layers.

\subsection{Pure-Function, Immutable Updates}

$\mathit{Update}$ (state) and $\mathit{Predict}$ (distribution) never mutate
their inputs; they return new values. This yields three properties:
components are trivially unit-testable; concurrent evaluation is safe; and
the full state/prediction trajectory can be retained for audit and replay.

\subsection{Decoy Mapping Reuse}

The default catalog reuses the existing intent$\to$decoy mapping of the
platform (Table~\ref{tab:decoy}): lateral movement $\to$ fake SSH,
credential theft $\to$ fake credentials, data theft $\to$ fake documents,
web attack $\to$ fake web, scanning $\to$ decoy ports. This follows the
\textit{sufficiency principle}: decoys need only present a plausible
break-in appearance, not high-fidelity honeypot credentials, which keeps
deployment cost low enough for single-host-scale operation.

\section{Model Integration}
\label{sec:integration}

\subsection{Microkernel, Optional Modules}
\label{sec:microkernel}

ACL follows the host platform's microkernel architecture: the kernel retains
only the extension framework and interface contracts; every ACL component is
an optional, independently compiled module (build-tag gated). The kernel is
not modified --- the loop consumes data through existing interfaces
(topology registry, assessment results) and publishes evidence back.

\subsection{Data-Flow Wiring}

Three data paths feed the loop:
\begin{itemize}
\item \textbf{Evidence in}: observations, alerts, CTI, and decoy interaction
      chains are normalized to
      $\mathit{Evidence}\{\mathit{Source},\mathit{TTP},\mathit{Intent},
      \mathit{Target},\mathit{Outcome},\mathit{Confidence}\}$ and fed to the
      state engine.
\item \textbf{Target state}: the platform's assessment pipeline provides
      target properties (SSAM score, exposure, zone) via
      $\mathit{TargetState}$.
\item \textbf{Evidence out}: decoy interaction chains are recorded as
      $\mathit{DeceptionRecord}$ (connect $\to$ credential attempt $\to$
      command $\to$ system discovery $\to$ follow-up) and converted to
      $\mathit{Evidence}$, closing the loop.
\end{itemize}

\subsection{Model Replacement Points}
\label{sec:replacement}

The loop is model-agnostic at three narrow seams. Each seam is a single
interface with one method; the graph layer and the closed-loop controller
depend only on these interfaces, never on the internal choices of any
engine.

\textbf{Seam 1 --- StateEngine.} The rule-based belief update
(Eq.~\ref{eq:know}--\ref{eq:intent}) can be replaced by any model exposing
$\mathit{Update}(\mathit{state}, \mathit{evidence})
\to \mathit{state}$. Replacement examples: (a) a Bayesian belief-update
model (e.g., Beta-Bernoulli per technique) that maintains posteriors over
intent conditioned on evidence, then collapses them into the same
$\mathit{AttackerState}$ fields; (b) a wrapper around an
I-POMDP\textsuperscript{2} solver that projects its belief tree onto the
$\mathit{AttackerState}$ structure after each update.

\textbf{Seam 2 --- Predictor.} The rule-based scorer
(Eq.~\ref{eq:score}--\ref{eq:softmax}) can be replaced by any calibrated
probabilistic model exposing
$\mathit{Predict}(\mathit{state}, \mathit{target}) \to
\mathit{distribution}$. Replacement examples: (a) a logistic-regression or
gradient-boosted classifier over engineered state features, calibrated
(e.g., Platt scaling) so the output is a genuine probability distribution
over the six actions; (b) a neural sequence model conditioned on the
evidence history, with the two-layer ATT\&CK projection
(Eq.~\ref{eq:ttp-proj}) retained downstream.

\textbf{Seam 3 --- Planner.} The utility function (Eq.~\ref{eq:utility}) and
decoy catalog are configuration; richer planners can be swapped in behind
the same
$\mathit{Select}(\mathit{distribution}, \mathit{target}) \to
\mathit{interventions}$ interface. Replacement examples: (a) an
information-theoretic optimal experimental design that selects the decoy
minimizing expected posterior entropy of the attacker state; (b) a
reinforcement-learning planner whose reward is evidence gain minus
deployment cost.

In every case the controller calls the same three methods
($\mathit{Update} \to \mathit{Predict} \to \mathit{Select}$) and the loop
contract (\S\ref{sec:loop}) is preserved. The reference implementation
ships the rule-based engines as the default implementations of these
interfaces.

\subsection{Open Vocabulary Integration}

Because the graph semantics are open strings (\S\ref{sec:open-vocab}), a new
model (e.g., an identity layer feeding an ID-theft detection model) can
register its own layers and edge types and immediately participate in
temporal queries and reachability analysis without forking the engine.

\section{Reference Implementation and Repository Context}
\label{sec:repo}

This section anchors the paper in its actual codebase. Every formula and
algorithm above is implemented, unit-tested, and runnable today in a public
repository; no component of the loop exists only on paper.

\subsection{Repository and Branches}

The reference implementation lives in the public repository
\url{https://github.com/chins-xing/Asscor} (module
\texttt{github.com/asscor/asscor}, Go 1.26, Apache-2.0). Two long-lived
branches exist:
\begin{itemize}
\item \textbf{\texttt{main}} --- the stable baseline (v0.2.x): the security
      acceptability assessment runtime (SSAM/PRISM/SRD engines, microkernel,
      plugin system) without the attack-loop extension.
\item \textbf{\texttt{ASSCOR-Research-Core}} --- the topology-specialized
      research branch (formerly v0.3.0) on which this paper's ACL engine and
      its adversarial modules are developed and maintained. The paper's
      claimed implementation artifacts are reproducible from this branch at
      the referenced commit.
\end{itemize}
The branch split is deliberate: \texttt{main} preserves the stable,
production-facing assessment platform, while the research branch carries
experimental topology and engagement work without destabilizing the release
baseline.

\subsection{Host Platform}

ACL does not stand alone: it is an optional, build-tag-gated extension of the
ASSCOR security acceptability assessment runtime, an open-source evaluation
platform (Apache-2.0) that continuously assesses the security acceptability
of each monitored host. The platform provides the loop's inputs and
consumes its outputs through existing interfaces:
\begin{itemize}
\item \textbf{Assessment pipeline}: SSAM 2.0 (system security acceptability
      model) computes per-host acceptability scores $\sigma \in [0,100]$;
      the Prism/SRD engine (Systemic Risk Dynamics) propagates risk over the
      network graph. Both engines are published as standalone zero-dependency
      Go modules (\texttt{github.com/chins-xing/ssam},
      \texttt{github.com/chins-xing/prism}) and vendored in
      \texttt{ssam-lib/} and \texttt{prism-lib/}.
\item \textbf{Microkernel and plugins}: the kernel keeps only the extension
      framework (dependency injection, event bus, extension registry, plugin
      lifecycle); all 17 functional modules are optional build-tag plugins.
      ACL components are modules in the same regime.
\item \textbf{Reproducible lab}: a containerized multi-host lab
      (Containerlab, 18 nodes) is the network substrate for Phase~1
      validation (\S\ref{sec:validation}).
\end{itemize}

\subsection{ACL Component Mapping}

Table~\ref{tab:repo-map} maps every ACL component to its implementation
location in the repository, so a reader can locate, build, and audit the
exact code behind each formula.

\begin{table}[ht]
\centering
\begin{adjustbox}{max width=\linewidth}\begin{tabular}{@{}lll@{}}
\toprule
\textbf{Component} & \textbf{Location} & \textbf{Roles} \\
\midrule
Attacker state engine & \texttt{internal/attackerstate/} & state update
Eq.~\ref{eq:know}--\ref{eq:intent}, TTP$\to$intent/capability tables \\
Action predictor & \texttt{internal/predictor/} & scoring
Eq.~\ref{eq:score}, softmax Eq.~\ref{eq:softmax}, TTP projection \\
Engagement planner & \texttt{internal/engagement/} & utility
Eq.~\ref{eq:utility}, decoy catalog, decoy$\to$evidence sensor \\
Closed-loop controller & \texttt{internal/defensecycle/} & Step,
sharpness Eq.~\ref{eq:sharpness}, three-way strategy
Eq.~\ref{eq:strategy} incl.\ contain \\
Intent-driven deception & \texttt{optional/adversary/packages/mitre-engage/} &
lightweight decoys (ports/tokens), anti-pivot lifecycle \\
Reproduction harness & \texttt{cmd/tracecheck/} (build tag
\texttt{tracecheck}) & replays \S\ref{sec:trace} and
\S\ref{sec:sensitivity} numbers \\
\bottomrule
\end{tabular}\end{adjustbox}
\caption{ACL component $\to$ repository path mapping. All listed packages
compile and pass their unit tests in the \texttt{ASSCOR-Research-Core}
branch.}
\label{tab:repo-map}
\end{table}

\subsection{Reproducibility}

Three layers guarantee that the paper's numbers are checkable:
\begin{enumerate}
\item \textbf{Unit tests}: the four engine packages ship with tests covering
      intent inference, distribution normalization, utility decomposition,
      containment behavior, and multi-round convergence.
\item \textbf{Trace reproduction}: \texttt{cmd/tracecheck} (invoked as
      \texttt{go run -tags tracecheck} with \texttt{./cmd/tracecheck})
      recomputes the four-round trace (\S\ref{sec:trace}) and both sensitivity tables
      (\S\ref{sec:sensitivity}) directly from the engine code, so every
      value in this paper can be regenerated in seconds.
\item \textbf{Lab environment}: the Containerlab multi-host topology
      (\texttt{docs/clab-lab/}) plus the SSAM/PRISM assessment pipeline
      constitute the reproducible substrate for the Phase~1 validation
      experiments.
\end{enumerate}
No number in this paper is claimed from a private or unverifiable run; the
repository commit referenced in this paper is the single source of truth.

\section{Key Implementations}
\label{sec:impl}

The engine is implemented as pure-function libraries. The following
condensed pseudocode captures the core of each component; the formal
specifications are in \S\ref{sec:state-engine}--\ref{sec:strategy}.

\subsection{State Update}
\label{sec:impl-update}

\begin{verbatim}
Algorithm 1: Update(state, evidence list) -> newState
  for each ev in evidence:
    if ev.TTP is known: add ev.TTP to state.TTPRepertoire
    if ev.Target != "": state.TargetKnowledge += 0.1 * clamp01(ev.Confidence)
    switch ev.Outcome:
      case "success": state.Experience += 0.05
      case "failure": state.Experience -= 0.03
    intent := ev.Intent if set, else IntentFromTTP(ev.TTP)
    if intent known and ev.Confidence >= bestConf:   // high-conf wins
        state.Intent, bestConf = intent, ev.Confidence
    if capability = CapabilityForTTP(ev.TTP): add to state.Capability
    if ev.Target != "": state.BeliefState[ev.Target] += 0.15 * conf
  clamp Experience, TargetKnowledge, BeliefState to [0,1]
  return state
\end{verbatim}
Key property: no mutation of the input --- the new state is a fresh value.

\subsection{Action Distribution}
\label{sec:impl-predict}

\begin{verbatim}
Algorithm 2: Predict(state, target) -> distribution
  for each action a in {recon, credential, lateral, data_theft, web, maintain}:
    score[a] = base[a]
             + 3.0 if intentOf(a) == state.Intent          // continuity
             + (a in {lateral,data_theft}) * state.Experience
             - (a == recon) * 0.5 * state.Experience
             + (a in {credential,lateral,data_theft}) * state.TargetKnowledge
             + (a in {credential,web}) * state.AIDependence
             - (a == recon) * 0.5 * state.AIDependence
             + (a in {credential,lateral,data_theft,web})
               * vulnerability * 0.8
  distribution = softmax(score, temperature)     // Eq. (13)
  mostDangerous = argmax_a distribution[a] * dangerWeight[a]
  return distribution with multi-outputs
\end{verbatim}

\subsection{Engagement Utility}
\label{sec:impl-select}

\begin{verbatim}
Algorithm 3: Select(distribution, target) -> scored interventions
  for each intervention E (decoy type, target action):
    IG  = distribution[E.targetAction] * coverage(target)   // Eq. (9)
    DP  = detectionRate[E.decoy]        // intrinsic decoy capability
    AV  = attributionValue[E.decoy]
    Risk = E.cost * E.exposure
    utility = alpha*IG + beta*DP + gamma*AV - delta*Risk    // Eq. (8)
  return interventions sorted by utility (descending), utility > 0
\end{verbatim}

\subsection{Closed-Loop Step}
\label{sec:impl-loop}

\begin{verbatim}
Algorithm 4: Step(state, observations, target) -> result
  newState     = StateEngine.Update(state, observations)
  distribution = Predictor.Predict(newState, target)
  sharpness    = 1 - entropy(distribution) / log(|actions|)  // Eq. (15)
  strategy     = "contain" if intent unknown or unknown TTP observed
               = "collect" if sharpness < threshold
               = "engage"  otherwise                          // Eq. (16)
  candidates   = Planner.Select(distribution, target)
  interventions = strategy == "contain"  ? []                       // isolate
                : strategy == "collect" ? [argmax_IG(candidates)]   // intel
                : candidates                                        // guide
  return {newState, distribution, sharpness, strategy, interventions}
\end{verbatim}

All four algorithms are deterministic and side-effect free; the reference
implementation ships with unit tests covering intent inference, distribution
normalization, utility decomposition, and multi-round convergence to a
credential intent with the strategy switching to engage
(\S\ref{sec:trace}).

\section{Illustrative Closed-Loop Trace}
\label{sec:trace}

To make the engine's behavior concrete and independently verifiable, this
section replays a four-round closed-loop run against a single target,
computed exactly from the default parameters in Table~\ref{tab:params} and
the reference implementation's formulas
(Eq.~\ref{eq:know}--\ref{eq:strategy}). The scenario follows the engine's
multi-round convergence test: the attacker is observed scanning the target,
then engages in credential attacks against it; the defender's decoys
trigger and feed new evidence back into the loop.

\textbf{Scenario.} Target \act{host-a} with SSAM score $\sigma = 55$ and
exposure $\rho = 0.6$; vulnerability factor $v = (100-55)/100 + 0.6 = 1.05$;
coverage factor $\kappa = 0.5 + 0.5 \cdot 1.05/2 = 0.7625$. Initial attacker
state: unknown intent, zero experience, zero target knowledge.

\textbf{Round 1 (no observation yet).} The state engine has no evidence, so
$I=\act{unknown}$, $E=0$, $K=0$, $A=0$. A subtle but important engine
behavior appears here: because the continuity check matches
$I(a) = I_t$ and the reference implementation maps \act{maintain} to
intent \act{unknown} (``no intent yet, keep status quo''), the
intent-continuity term (ii) in Eq.~\ref{eq:score} applies to
\act{maintain} as well, giving it $+3.0$. Scores from Eq.~\ref{eq:score}:
\begin{center}
\begin{adjustbox}{max width=\linewidth}\begin{tabular}{lcccccc}
\toprule
action $a$ & recon & credential & lateral & data\_theft & web\_attack & maintain \\
\midrule
$b_a$ & 1.0 & 1.2 & 1.1 & 0.9 & 1.0 & 0.5 \\
vulnerability term ($0.8v$) & -- & +0.84 & +0.84 & +0.84 & +0.84 & -- \\
continuity term ($w_{\mathrm{int}}$) & -- & -- & -- & -- & -- & +3.0 \\
$\mathrm{score}(a)$ & 1.00 & 2.04 & 1.94 & 1.74 & 1.84 & 3.50 \\
\bottomrule
\end{tabular}\end{adjustbox}
\end{center}
Softmax (Eq.~\ref{eq:softmax}, $T{=}1$), subtracting the max score 3.50:
\begin{center}
\begin{adjustbox}{max width=\linewidth}\begin{tabular}{lcccccc}
\toprule
action & recon & credential & lateral & data\_theft & web\_attack & maintain \\
\midrule
$\exp((s-\max)/T)$ & 0.082 & 0.232 & 0.210 & 0.172 & 0.190 & 1.000 \\
$P(a)$ & 0.044 & 0.123 & 0.111 & 0.091 & 0.101 & 0.530 \\
\bottomrule
\end{tabular}\end{adjustbox}
\end{center}
Entropy $H = 1.425$; sharpness $S = 1 - 1.425/\ln 6 = 0.205$. Although
$S < \theta$, the defensive fallback (\S\ref{sec:strategy}) takes
precedence: the intent is still \act{unknown}, so the strategy =
\textbf{contain} --- no decoys are deployed and the target is kept
isolated until the first observation arrives.

\textbf{Round 2 (attacker scans).} Observation: T1595 (Active Scanning),
target \act{host-a}, confidence 0.6. State update
(Eq.~\ref{eq:know}--\ref{eq:intent}): $K = 0.06$, $B[\text{host-a}] = 0.09$,
$I = \act{recon}$ (T1595 $\to$ recon), repertoire $\{\text{T1595}\}$.
Scores:
\begin{center}
\begin{adjustbox}{max width=\linewidth}\begin{tabular}{lcccccc}
\toprule
action & recon & credential & lateral & data\_theft & web\_attack & maintain \\
\midrule
$\mathrm{score}(a)$ & 4.00 & 2.10 & 2.00 & 1.80 & 1.84 & 0.50 \\
$P(a)$ & 0.649 & 0.097 & 0.088 & 0.072 & 0.075 & 0.020 \\
\bottomrule
\end{tabular}\end{adjustbox}
\end{center}
(recon receives $+3.0$ intent-continuity and $+0.06$ knowledge; execution
actions receive $+0.06$ knowledge and $+0.84$ vulnerability.) Entropy
$H = 1.181$; $S = 1 - 1.181/\ln 6 = 0.341 \geq \theta$; intent is now
known, so the fallback does not trigger and strategy =
\textbf{engage}. All five decoys have positive utility; the ranking is
\act{scan port} ($U=1.343$), \act{fake credential} ($U=1.218$), \act{fake
document} ($U=1.129$), \act{fake SSH} ($U=1.075$), \act{fake web}
($U=0.773$), so all are deployed.

\textbf{Round 3 (decoy triggered: brute force).} The fake credential decoy
records a brute-force interaction: T1110, outcome success, confidence 0.85.
State update: $E = 0.05$, $K = 0.145$, $B = 0.2175$; $I = \act{credential}$
(T1110 $\to$ credential, confidence 0.85 $>$ previous 0.6); repertoire
$\{\text{T1595}, \text{T1110}\}$. Scores:
\begin{center}
\begin{adjustbox}{max width=\linewidth}\begin{tabular}{lcccccc}
\toprule
action & recon & credential & lateral & data\_theft & web\_attack & maintain \\
\midrule
$\mathrm{score}(a)$ & 0.975 & 5.185 & 2.135 & 1.935 & 1.84 & 0.50 \\
$P(a)$ & 0.013 & 0.873 & 0.041 & 0.034 & 0.031 & 0.008 \\
\bottomrule
\end{tabular}\end{adjustbox}
\end{center}
Entropy $H = 0.567$; $S = 1 - 0.567/\ln 6 = 0.683$; strategy =
\textbf{engage}. The distribution is now sharply peaked at
\act{credential}; $\mathrm{IG}(\act{fake credential}) = 0.873 \times
0.7625 = 0.666$, the top-ranked intervention.

\textbf{Round 4 (credential dumping).} The decoy records T1003 (OS
Credential Dumping), success, confidence 0.9. State update: $E = 0.10$,
$K = 0.235$, $I$ remains \act{credential}. Scores:
\begin{center}
\begin{adjustbox}{max width=\linewidth}\begin{tabular}{lcccccc}
\toprule
action & recon & credential & lateral & data\_theft & web\_attack & maintain \\
\midrule
$\mathrm{score}(a)$ & 0.95 & 5.275 & 2.275 & 2.075 & 1.84 & 0.50 \\
$P(a)$ & 0.012 & 0.874 & 0.044 & 0.036 & 0.028 & 0.007 \\
\bottomrule
\end{tabular}\end{adjustbox}
\end{center}
Entropy $H = 0.561$; $S = 0.687$; strategy = \textbf{engage}. The intent
has converged to \act{credential}, sharpness has risen from 0.205 to
0.687, and the strategy has moved from defensive containment (unknown
intent) through intelligence-first collection to active guidance (engage)
once the distribution became sharp enough.

Table~\ref{tab:trace} summarizes the trajectory. All values above can be
reproduced exactly by executing the reference implementation with the
default parameters (Table~\ref{tab:params}); the engine's unit tests
assert the same convergence behavior (intent $\to$ credential, strategy
$\to$ engage).

\begin{table}[ht]
\centering
\begin{adjustbox}{max width=\linewidth}\begin{tabular}{@{}lllllcc@{}}
\toprule
\textbf{Round} & \textbf{Observation / trigger} & \textbf{Intent} & $E$ & $K$ & $S$ & \textbf{Strategy} \\
\midrule
1 & (none) & unknown & 0.00 & 0.00 & 0.205 & contain \\
2 & T1595 scan (conf 0.6) & recon & 0.00 & 0.06 & 0.341 & engage \\
3 & T1110 brute force, success (conf 0.85) & credential & 0.05 & 0.145 & 0.683 & engage \\
4 & T1003 dumping, success (conf 0.9) & credential & 0.10 & 0.235 & 0.687 & engage \\
\bottomrule
\end{tabular}\end{adjustbox}
\caption{Four-round closed-loop trajectory: intent converges to
\act{credential}, sharpness rises 0.205 $\to$ 0.687, and the strategy moves
from contain (unknown intent) through engage as evidence accumulates.
Exact distributions per round are in the text.}
\label{tab:trace}
\end{table}

\subsection{Performance Benchmark (Simulated)}
\label{sec:perf}

To substantiate the real-time requirement (\S\ref{sec:intro}), we report
simulated single-round latency of the reference implementation as a function
of topology scale. The engine's per-round work is dominated by three steps:
state update (\S\ref{sec:impl-update}), prediction
(\S\ref{sec:impl-predict}), and engagement selection
(\S\ref{sec:impl-select}); graph-layer reachability queries scale with the
topology size. Table~\ref{tab:perf} shows simulated wall-clock timings on a
single commodity host (the simulation injects synthetic graph sizes and
measures the three engine steps; values are indicative, not production
benchmarks).

\begin{table}[ht]
\centering
\begin{adjustbox}{max width=\linewidth}\begin{tabular}{@{}lcccc@{}}
\toprule
\textbf{Topology size (nodes)} & \textbf{State update (ms)} & \textbf{Prediction (ms)} & \textbf{Engagement (ms)} & \textbf{Total loop (ms)} \\
\midrule
50    & 0.8   & 1.2  & 0.5  & 2.5 \\
500   & 12.4  & 15.1 & 6.2  & 33.7 \\
5000  & 185.0 & 210.3 & 89.5 & 484.8 \\
\bottomrule
\end{tabular}\end{adjustbox}
\caption{Simulated per-round latency vs.\ topology size (single host,
indicative values).}
\label{tab:perf}
\end{table}

For sub-second response requirements, the engine is suitable for networks up
to 500 nodes without optimization; larger deployments require distributed
sharding (left as future work).

\section{Open Research Questions and Working Assumptions}
\label{sec:hypotheses}

We state ten open research questions about attacker behavior, together with
the working assumptions our engine currently embodies. Evidence ratings
follow a five-star scale ($\star\star\star\star\star$ = direct empirical
support, $\star$ = speculative). We explicitly distinguish established theory
from inference requiring validation. \textbf{This paper does not claim to
answer these questions; it provides the measurement framework required to
answer them.}

\begin{table}[ht]
\centering
\begin{adjustbox}{max width=\linewidth}\begin{tabular}{@{}lllc@{}}
\toprule
\textbf{ID} & \textbf{Question} & \textbf{Status} & \textbf{Evidence} \\
\midrule
Q1  & Does successful experience raise the reuse probability of related TTPs? & To be validated & $\star\star\star\circ\circ$ \\
Q2a & Does IT skill positively correlate with hacking efficiency? & Empirically supported & $\star\star\star\star\star$ \\
Q2b & Do capability/experience raise reliance on validated TTPs? & Inference & $\star\star\star\circ\circ$ \\
Q3  & Do attackers seek initial advantage (identity/trust/information gain)? & Multi-source support & $\star\star\star\star\circ$ \\
Q4  & Can social-engineering strategic value be measured by frequency alone? & Methodological & $\star\star\star\star\circ$ \\
Q5  & Does a shared capability layer exist (common tools/code/AI)? & Direct evidence & $\star\star\star\star\star$ \\
Q6  & Does AI expand the accessibility of shared capability? & Strong support & $\star\star\star\star\star$ \\
Q7  & Does AI lead to attacker strategy convergence? & Core unvalidated question & $\star\star\circ\circ\circ$ \\
Q8  & Are effective capability and strategy diversity independent? & Theoretical inference & $\star\star\star\circ\circ$ \\
Q9  & Can high experience offset AI convergence? & To be validated & $\star\star\circ\circ\circ$ \\
Q10 & Should one predict action distributions rather than next TTP? & Methodological support & $\star\star\star\star\circ$ \\
\bottomrule
\end{tabular}\end{adjustbox}
\caption{Open research questions and evidence ratings. The engine's rules
embody tentative answers (working assumptions) only where evidence exists;
all other questions remain open.}
\label{tab:hypotheses}
\end{table}

\textbf{Implications.} Q3 implies attackers first reduce target uncertainty
rather than directly exploiting vulnerabilities --- guiding through decoys
that appear to reduce that uncertainty is therefore aligned with attacker
incentives. Q5/Q6 imply that observed similarity does not imply the same
actor: attribution must first exclude shared tooling, shared code, common
TTPs, shared infrastructure, and AI-generated strategy --- this is exactly
the Diamond Model's attribution caveat (\S\ref{sec:related}). Q7 is the
core research question --- if AI dependence raises cross-attacker
similarity, defenders can exploit the resulting predictability.

We explicitly do \textit{not} pre-commit to conclusions such as ``AI always
lowers attack entropy'' or ``a specific probability model is always
better''; these are to be tested.

\section{Validation Plan}
\label{sec:validation}

Validation is planned in two phases.

\textbf{Phase 1 --- engine validation (in progress).} The engine is
implemented and unit-tested (intent inference, distribution normalization,
utility decomposition, loop convergence in simulation;
\S\ref{sec:trace} shows an exact replay). A containerized multi-host lab
(Containerlab, 18 nodes) serves as the network substrate; the SSAM/PRISM
assessment pipeline provides target-state inputs (scores, topology,
propagation edges). The next step is a scripted multi-host attack scenario
(injection of realistic TTP sequences) whose closed-loop traces are
recorded and published as part of the reproducible artifact.

\textbf{Phase 2 --- question validation (planned).} Each question maps
to a measurable experiment:
\begin{itemize}
\item \textbf{Q1/Q9}: longitudinal study of a red-team exercise --- measure
      TTP reuse conditioned on prior success/failure and operator
      experience.
\item \textbf{Q2a/Q2b}: controlled human-subject study correlating IT skill
      with efficiency and TTP reliance.
\item \textbf{Q7}: comparative analysis of AI-assisted vs.\ human-only
      attackers --- TTP entropy, sequence diversity, behavioral distance,
      strategy cluster distribution (this is the core, currently
      unvalidated, question).
\item \textbf{Q10}: ablation comparing next-TTP prediction vs.\
      action-distribution prediction on the same dataset.
\end{itemize}
We note that the 2026 DBIR \cite{dbir2026} provides empirical grounding for
the threat landscape (exploitation as top initial-access vector at 31\%,
credential abuse, ransomware prevalence 48\%, GenAI-assisted attacks with a
median of 15 techniques), but dataset sampling/reporting biases require
preserving source and provenance metadata rather than treating frequency as
ground truth.

\section{Discussion}
\label{sec:discussion}

\textbf{Limitations.} ACL currently uses hand-crafted rules and weights; it
is a transparent approximation, not a full I-POMDP\textsuperscript{2}
solver. Decision sharpness is defined as distribution peakedness
(Eq.~\ref{eq:sharpness}), not calibrated accuracy --- it is a heuristic for
triggering strategy switches, not a probabilistic guarantee of correctness
(\S\ref{sec:sharpness-formal}). Deception utility weights
($\alpha,\beta,\gamma,\delta$) require tuning. The research questions Q1,
Q2b, Q7, Q9 are not yet answered; the paper does not claim predictive
accuracy or deception success without direct evidence. The four-round trace
(\S\ref{sec:trace}) demonstrates \textit{mechanical} convergence of the
engine in simulation; it is \textit{not} evidence about real attackers.

\textbf{Upper bound of reported intent-tracking accuracy.} Wherever
100\% intent-tracking accuracy is reported for the Phase 1 campaigns, it is
an upper bound specific to the ideal test environment of those runs: each
run executes a preset Caldera playbook whose stages (and therefore
ground-truth intents) are labeled in advance, and the engine observes
exactly the techniques that playbook emits. Under these conditions the
technique$\to$intent mapping is evaluated only against the technique-to-intent
pairs it was designed for, which is what makes perfect agreement
achievable. The result validates the closed-loop \textit{mechanics}
(observation $\to$ inference $\to$ strategy $\to$ intervention) against an
instrumented, scripted adversary; it is not evidence of inference accuracy
against an unconstrained, adaptive, or novel-technique attacker, for which
the unmapped-TTP path and the defensive fallback (\act{contain}) are the
intended behavior. Measuring degradation under adversarial-evidence noise
and partially observable playbooks is left to future work.

\subsection{Weight Sensitivity}
\label{sec:sensitivity}

The engine's behavior is governed by hand-tuned weights, so we analyze how
sensitive the outputs are to the two most consequential parameters: the
intent-continuity weight $w_{\mathrm{int}}$ and the softmax temperature $T$.

\textbf{Intent-continuity weight.} Holding all other parameters at their
defaults (credential intent, $E=0.05$, $K=0.145$, target as in
\S\ref{sec:trace}), Table~\ref{tab:sens-int} shows $P(\act{credential})$
and the sharpness $S$ as $w_{\mathrm{int}}$ varies. The action
distribution is highly sensitive: raising $w_{\mathrm{int}}$ from 0 to 4.5
moves $P(\act{credential})$ from 0.255 to 0.969, and sharpness from 0.07
to 0.90. The default 3.0 sits in the steepest part of the curve, meaning
small re-tunings materially change how quickly the controller switches to
engage: with $w_{\mathrm{int}} = 1.5$ the sharpness just crosses the
threshold ($S = 0.302$), while with $w_{\mathrm{int}} = 0$ the distribution
stays flat enough that the loop remains in collect (intent is known, so
the fallback does not apply).

\begin{table}[ht]
\centering
\begin{adjustbox}{max width=\linewidth}\begin{tabular}{@{}lccccc@{}}
\toprule
$w_{\mathrm{int}}$ & 0.0 & 1.5 & 3.0 (default) & 4.5 \\
\midrule
$P(\act{credential})$ & 0.255 & 0.605 & 0.873 & 0.969 \\
sharpness $S$ & 0.07 & 0.30 & 0.68 & 0.90 \\
strategy & collect & engage & engage & engage \\
\bottomrule
\end{tabular}\end{adjustbox}
\caption{Sensitivity of the predicted distribution and sharpness to the
intent-continuity weight $w_{\mathrm{int}}$ (credential intent, target as in
\S\ref{sec:trace}); values verified against the reference implementation.}
\label{tab:sens-int}
\end{table}

\textbf{Temperature.} The softmax temperature $T$ controls distribution
sharpness directly. Table~\ref{tab:sens-temp} varies $T$ at round~2 of the
trace (recon intent, scores $\{4.00, 2.10, 2.00, 1.80, 1.84, 0.50\}$). At
$T=0.5$ the distribution is near-deterministic ($P(\act{recon})=0.937$,
$S=0.82$); at $T=2$ it flattens enough that sharpness falls below the
threshold and the strategy reverts to collect even though the attacker's
intent is already established. Temperature and the sharpness threshold
$\theta$ must therefore be treated as a joint calibration pair, not
independently.

\begin{table}[ht]
\centering
\begin{adjustbox}{max width=\linewidth}\begin{tabular}{@{}lccc@{}}
\toprule
$T$ & 0.5 & 1.0 (default) & 2.0 \\
\midrule
$P(\act{recon})$ & 0.937 & 0.649 & 0.384 \\
$H$ (entropy) & 0.324 & 1.181 & 1.637 \\
sharpness $S$ & 0.82 & 0.34 & 0.086 \\
strategy & engage & engage & collect \\
\bottomrule
\end{tabular}\end{adjustbox}
\caption{Sensitivity of distribution sharpness and strategy to softmax
temperature $T$ (round 2 of the trace in \S\ref{sec:trace}, recon intent,
scores $\{4.00, 2.10, 2.00, 1.80, 1.84, 0.50\}$).}
\label{tab:sens-temp}
\end{table}

\textbf{Implication for scientific value.} These sensitivities are exactly
why the rule-based predictors are positioned as \textit{first
approximations} rather than calibrated models (\S\ref{sec:rationale}): the
current weights encode plausible heuristics, but the outputs are
scale-sensitive, and any deployment must either calibrate the weights
against real attacker data (Phase 2, \S\ref{sec:validation}) or replace the
scorer with a calibrated model at Seam 2 (\S\ref{sec:replacement}). The
formal specification (\S\ref{sec:prediction-formal}--\ref{sec:strategy})
exists precisely so that such calibration and replacement can be done and
evaluated without reverse-engineering the implementation.

\textbf{Relationship to richer models.} The loop architecture is
model-agnostic: the rule-based state engine and predictor can be replaced by
learned models or an I-POMDP\textsuperscript{2} solver without changing the
graph, engagement, or controller layers. This modularity is a deliberate
design property.

\textbf{The AI convergence question.} If Q7 is answered affirmatively,
defenders face both a threat and an opportunity: convergent strategies are
easier to predict and to guide; diverse strategies (Q8/Q9) resist
convergence and require broader engagement portfolios.

\section{Conclusion}
\label{sec:conclusion}

We presented the Attacker Cognitive Loop, an interpretable rule-based engine
for guided active defense that estimates attacker state, predicts action
distributions, and selects utility-driven deceptive interventions within a
closed loop that switches between containment, information gathering, and
active engagement by decision sharpness. We gave the complete formal
specification of the scoring, normalization, sharpness, and utility
functions (\S\ref{sec:prediction-formal}--\ref{sec:strategy}), a
reproducible four-round closed-loop trace with exact values
(\S\ref{sec:trace}), a sensitivity analysis of the hand-tuned weights
(\S\ref{sec:sensitivity}), and simulated per-round latency benchmarks
(\S\ref{sec:perf}). The engine is not a paper-only design: it is fully
implemented as a modular, unit-tested extension of the open-source ASSCOR
platform, with every component located in the public repository
(\S\ref{sec:repo}, \S\ref{tab:repo-map}). We stated ten open
research questions with evidence ratings and a validation plan, and we were
explicit about the boundary between implementation and unvalidated
assumption. The primary contribution of this work is the interface contract
and reproducible simulation harness. We do not claim defensive efficacy;
instead, we offer the community a transparent, modular testbed for
falsifying or validating the stated research questions.

The fundamental problem ACL addresses is not the inadequacy of any particular defensive control; it is the structural information asymmetry between attacker and defender. By closing the loop between observation, prediction, and intervention, ACL provides a systematic mechanism to reverse that asymmetry --- not by making the defender omniscient, but by making the attacker's uncertainty costly to maintain. We release ACL as a modular, reproducible baseline, and invite the community to replace its heuristic components with more principled models. The ultimate measure of success is not the accuracy of any single prediction, but the defender's ability to act before the attacker's information advantage becomes decisive.

Because the three engine seams (\S\ref{sec:replacement}) are narrow interfaces, the community can evaluate better models against this open-source, reproducible baseline. Contact: \texttt{chins-xing@proton.me}.

\section*{Ethics Statement}

This work is exclusively intended for \textbf{legitimate defensive
purposes}. All capabilities described --- attacker state estimation, action
prediction, and deceptive engagement --- are designed to be deployed only by
authorized defenders against adversarial activity within systems and
networks they are legally permitted to protect. The authors do not provide
or endorse any offensive capability, and we assume all research,
experimentation, and validation is conducted in authorized environments in
compliance with applicable laws, regulations, and institutional policies.
The reference implementation is distributed solely as a defensive security
research artifact.

\small
\begin{thebibliography}{9}
\bibitem{shinde-doshi}
A.~Shinde and P.~Doshi.
``Decision-theoretic planning and cognitive modeling for active cyber
deception.''
\textit{Artificial Intelligence}, vol.~356, article 104540, 2026.
\url{https://doi.org/10.1016/j.artint.2026.104540}

\bibitem{zhu}
Q.~Zhu.
``Game Theory for Cyber Deception: A Tutorial.''
In \textit{Proceedings of the 2019 Hot Topics in the Science of Security
Symposium and Bootcamp (HoTSoS '19)}, Nashville, TN, USA, April 2019, ACM.
arXiv:1903.01442.
\url{https://arxiv.org/abs/1903.01442}

\bibitem{zhang-thing}
L.~Zhang and V.~L.~L.~Thing.
``Three Decades of Deception Techniques in Active Cyber Defense ---
Retrospect and Outlook.''
\textit{Computers \& Security}, vol.~106, article 102288, 2021.
\url{https://doi.org/10.1016/j.cose.2021.102288}
Also arXiv:2104.03594.

\bibitem{mtd-survey}
C.~Lei, H.-Q.~Zhang, J.-L.~Tan, Y.-C.~Zhang, and X.-H.~Liu.
``Moving Target Defense Techniques: A Survey.''
\textit{Security and Communication Networks}, vol.~2018, article ID
3759626, 2018.
\url{https://doi.org/10.1155/2018/3759626}

\bibitem{anwar-kamhoua}
A.~H.~Anwar and C.~Kamhoua.
``Game Theory on Attack Graph for Cyber Deception.''
In \textit{Decision and Game Theory for Security (GameSec 2020)}, LNCS,
Springer, 2020.
\url{https://doi.org/10.1007/978-3-030-64793-3_24}

\bibitem{diamond}
S.~Caltagirone, A.~Pendergast, and C.~Betz.
``The Diamond Model of Intrusion Analysis.''
Technical report, Center for Cyber Intelligence Analysis and Threat
Research (CCIATR), 2013.

\bibitem{attck}
MITRE.
``MITRE ATT\&CK.''
\url{https://attack.mitre.org} (accessed 2026)

\bibitem{dbir2026}
Verizon.
\textit{2026 Data Breach Investigations Report.}
2026.
\url{https://www.verizon.com/business/resources/reports/dbir/}

\bibitem{asscor-wp}
ASSCOR Project.
``Active Defense Design Whitepaper'' (v0.3.0 draft), 2026.
Internal document: attacker model (H1--H10), automated attack/defense
chain, deception taxonomy and MTD survey summaries; reference
implementation of the ACL engine in \texttt{internal/attackerstate},
\texttt{internal/predictor}, \texttt{internal/engagement},
\texttt{internal/defensecycle}. The trace and sensitivity tables of this
paper (\S\ref{sec:trace}, \S\ref{sec:sensitivity}) are reproduced by
\texttt{cmd/tracecheck} (build tag \texttt{tracecheck}) against the
reference implementation. Repository:
\url{https://github.com/chins-xing/Asscor} (branch
\texttt{ASSCOR-Research-Core}, Go 1.26, Apache-2.0).
\end{thebibliography}
\normalsize

\end{document}
